\label{sec:conclusion}

We have documented a strong negative correlation ($r = -0.67$) between district
Polsby-Popper compactness and mean constituent driving time across 432 congressional
districts, a 14-minute advantage in mean driving time for \textsc{bisect} algorithmic
plans relative to enacted gerrymanders in the 12-state partisan gerrymandering subsample,
and a geometric mechanism---the district neck---that explains why compact maps produce
shorter constituent travel distances.

\subsection{Implications for Redistricting Law}

The constituent distance finding has a straightforward legal implication: compactness
requirements, properly enforced, protect constituents from a concrete, welfare-relevant
harm. Courts that have hesitated to operationalize compactness requirements (beyond
prohibiting ``bizarre'' shapes) now have empirical content for the requirement: a
compact district is one that does not impose excessive constituent travel burdens.

This implication is non-partisan. The 14-minute advantage in mean driving time under
\textsc{bisect} plans benefits constituents in Democratic and Republican districts
equally---what matters is the geometry, not the partisan composition. In states where
the gerrymander benefits Democrats (Maryland), Democratic constituents in non-compact
districts face the same elevated travel times as Republican constituents in non-compact
districts in states where the gerrymander benefits Republicans (North Carolina, Ohio,
Texas). The harm is geometric, not partisan.

\subsection{Limitations}

Several limitations qualify the analysis. First, we use the population centroid as the
representative office location for \textsc{bisect} plans, which understates the
true benefit if actual representatives under algorithmic plans would choose office
locations optimally. Second, OSRM driving times do not account for traffic,
parking, or transit alternatives for urban constituents who may prefer transit.
Third, the analysis covers driving time only; access by mail, phone, and email
is not captured but is increasingly the primary constituent interaction modality.

\subsection{Future Work}

The O.4 analysis will be extended in two directions. First, we will apply the analysis
to \textsc{bisect} plans for all 50 states across all three census years (2000, 2010, 2020)
to provide a temporal view of whether enacted maps have diverged from compact algorithmic
alternatives in their constituent access properties. Second, we will analyze transit
travel times for urban districts using OSRM's public transit routing mode, providing
a more complete picture of access for urban constituents who do not primarily travel
by car.

\bigskip
\noindent\textit{The author thanks the OpenStreetMap contributors for the road network
data and the OSRM project for routing infrastructure. District office address data
was collected from the public \texttt{house.gov} website. All errors are the author's own.}
